\section{Auditoría de fuentes: reescribir la historia de crecimiento como mecanismos del sistema}

Esta clase la imparte Sualeh Asif, CTO y cofundador de Cursor. Al verificarlo el 11 de agosto de 2026, el video público ya era private; el repositorio conserva 1,288 subtítulos con marca de tiempo y la portada oficial. Los subtítulos registran la experiencia de arquitectura e incidentes de marzo de 2025; la documentación de Cursor de 2026 sobre secure indexing, privacy, CursorBench, Router y agent solo se usa para verificar mecanismos persistentes o explicar su evolución posterior, y no se proyecta hacia atrás como si esas capacidades ya existieran en el momento de la clase.

\begin{warningbox}{Límites de la evidencia sobre las cifras de escala}
En la clase se mencionan cifras como un crecimiento de aproximadamente cien veces en el último año, alrededor de cien millones de llamadas diarias a los modelos propios y alrededor de mil millones de documentos diarios en el sistema de indexación. Este texto las conserva como \textbf{estimaciones del ponente/cifras de la clase}, útiles para entender la forma de la carga, y no como conclusiones auditadas sobre la escala actual, la posición comercial o las relaciones con proveedores.
\end{warningbox}

\begin{importantbox}{Los tres ciclos cerrados de toda la clase}
\begin{enumerate}
  \item \textbf{Context loop}: workspace change $\rightarrow$ incremental index $\rightarrow$ retrieval context;
  \item \textbf{Edit loop}: request $\rightarrow$ model/provider $\rightarrow$ plan/edit $\rightarrow$ fast apply $\rightarrow$ validation;
  \item \textbf{Reliability loop}: traffic/failure $\rightarrow$ telemetry $\rightarrow$ mitigation/migration $\rightarrow$ safer architecture.
\end{enumerate}
\end{importantbox}

\subsection{Por qué AI coding no es una API de modelo}

La experiencia de usuario de AI coding depende en conjunto de la freshness del contexto, la retrieval quality, la capacidad del modelo, la provider capacity, la edit application y la editor validation. Aunque la salida del modelo sea correcta, si el índice está desactualizado, el provider tiene cola, la aplicación del patch es lenta o los diagnostics se traban, el usuario percibe que el sistema falla. Por eso la optimización de la infraestructura debe tomar como objeto la interaction loop completa.

\lecturefigure{01-coding-loop.png}{Un producto de AI coding es un ciclo cerrado formado por workspace, context, model, apply y feedback.}{Subtítulos locales 00:00--05:20; redibujo conceptual.}

\begin{knowledgebox}{Lectura de la figura: el latency budget pertenece a toda la cadena}
Una interacción puede descomponerse, de forma aproximada, como
\[
T_{total}=T_{context}+T_{queue}+T_{model}+T_{apply}+T_{validate}.
\]
El p95/p99 de cualquiera de estos términos basta para romper la fluidez. El streaming puede reducir el time-to-first-token, pero no elimina la latencia de cola de context lookup, queueing, patch conflict ni diagnostics.
\end{knowledgebox}

\teachervoice{Asif advierte a los estudiantes con la frase «si cien millones de llamadas no te suenan aterradoras, en realidad sí lo son»: el crecimiento no solo cambia la capacidad, sino cada supuesto por defecto. Los retry ocasionales, los rebuild completos o las reparaciones manuales que antes eran esporádicos se convierten, a gran cardinalidad, en carga continua.}

\subsection{Resumen de la sección}

El protagonista de Lecture 07 no es la marca Cursor, sino un AI interaction system de alta frecuencia. Cada decisión de arquitectura que sigue debe responder: qué tramo del loop mejora, qué estado agrega y cómo se degrada ante una falla.
